\begin{proof}[Proof of \cref{obj:thm:universal-point-identification}]
\begin{enumerate}
\item For \(a\in\{0,1\}\) and \(r\in\mathcal R_J\), set
\[
\delta_{ar}:=
\inf_{c\in[1/(1-\varepsilon),\,1/\varepsilon]}
\int_{C_r}|1-cp_a(e)|\,H(de),
\qquad p_1(e)=e,\quad p_0(e)=1-e.
\]
The coefficient interval is nonempty under overlap. Each defining integral is finite and nonnegative, so \(\delta_{ar}\geq0\). By \cref{obj:thm:all-label-ambiguity},
\[
D_H(g)=\sum_{r\in\mathcal R_J}(\delta_{1r}+\delta_{0r}).
\]
Thus \(D_H(g)=0\) implies \(\delta_{1r}=0\) for every \(r\): each term in this finite sum is nonnegative. % lean: CausalSmith.PartialID.UnlinkedPropensityAte.worstCaseAmbiguity_eq

\item Fix \(r\) with \(H(C_r)>0\), and define, for every real \(t\),
\[
\Phi_r(t):=\int_{C_r}|1-te|\,H(de).
\]
Since \(0\leq e\leq1\) on \(\mathcal E\), for all real \(t,t'\),
\[
|\Phi_r(t)-\Phi_r(t')|
\leq \int_{C_r}|t-t'|e\,H(de)
\leq |t-t'|H(C_r)
\leq |t-t'|.
\]
Hence \(\Phi_r\) is continuous and attains its minimum on the compact coefficient interval. Let \(c_r\) attain it. The preceding step gives
\[
0=\delta_{1r}=\Phi_r(c_r)
=\int_{C_r}|1-c_re|\,H(de),
\]
so \(1-c_re=0\) for \(H\)-almost every \(e\in C_r\). Because \(H(C_r)>0\), there is an \(x_r\in C_r\) with \(1-c_rx_r=0\). Also \(c_r\geq1/(1-\varepsilon)>0\). Consequently,
\[
c_r(e-x_r)=0
\quad\text{for \(H\)-almost every \(e\in C_r\),}
\]
and therefore \(e=x_r\) there almost surely. This proves the cellwise condition. % lean: CausalSmith.PartialID.UnlinkedPropensityAte.cellAbsoluteDeviation_minimizer

\item Conversely, suppose the cellwise condition holds. Choose its value \(x_r\) on each positive-mass cell and define the map on the entire label set by
\[
h(r):=
\begin{cases}
x_r,&H(C_r)>0,\\
\varepsilon,&H(C_r)=0,
\end{cases}
\qquad r\in\mathcal R_J.
\]
Here \(\varepsilon\in\mathcal E\) by overlap. The finite label set makes \(h\) measurable. On every positive-mass cell, \(e=h(r)\) almost surely; the same assertion is vacuous on a zero-mass cell. Taking the finite union of these cellwise assertions yields
\[
e=h(g(e))\qquad\text{for \(H\)-almost every \(e\)}.
\]
% lean: CausalSmith.PartialID.UnlinkedPropensityAte.scoreRecovery_iff_cellwiseAE

\item For every \(a\in\{0,1\}\) and \(r\in\mathcal R_J\), set
\[
c_{ar}:=\frac{1}{p_a(h(r))}.
\]
Overlap gives
\[
\varepsilon\leq p_a(h(r))\leq1-\varepsilon,
\qquad
\frac1{1-\varepsilon}\leq c_{ar}\leq\frac1\varepsilon.
\]
The cellwise equality established above gives \(p_a(e)=p_a(h(r))\) almost surely on \(C_r\), including when that cell has zero mass. Thus
\[
0\leq\delta_{ar}
\leq\int_{C_r}|1-c_{ar}p_a(e)|\,H(de)=0.
\]
Every summand in the all-label formula is zero, and hence \(D_H(g)=0\). Finally, if a measurable \(h\) satisfies \(e=h(g(e))\) \(H\)-almost surely, then on each positive-mass \(C_r\) its value \(h(r)\) witnesses the cellwise condition. The two stated conditions are therefore each equivalent to \(D_H(g)=0\). % lean: CausalSmith.PartialID.UnlinkedPropensityAte.worstCaseAmbiguity_eq_zero_iff
\end{enumerate}
\end{proof}
